\documentclass[12pt]{article}
% Préambule commun à tous les documents extraits de « La Revue CAPES »
\usepackage[T1]{fontenc}
\usepackage[utf8]{inputenc}
\usepackage{lmodern}
\usepackage{amsmath,amssymb,amsthm,amsfonts,mathrsfs}
\usepackage{setspace}
\usepackage{listings}
\usepackage{pgfplots}
\pgfplotsset{compat=1.18}
\usepackage{longtable}
\usepackage{adjustbox}
\usepackage{lettrine}
\usepackage{tkz-tab}
\usepackage{changepage}
\usepackage{pdflscape}
\usepackage{graphicx}
\usepackage{tikz}
\usepackage{xcolor}
\usepackage[a4paper,top=2cm,bottom=2.5cm,left=2.5cm,right=2cm]{geometry}
\usepackage{fancyhdr}
\usepackage{draftwatermark}
\SetWatermarkText{La RevueCapes}
\SetWatermarkLightness{0.9}
\SetWatermarkScale{2}
\usepackage{hyperref}
\hypersetup{colorlinks=true,linkcolor=blue!50!black,urlcolor=blue!50!black}

\definecolor{revue}{RGB}{20,60,120}

\pagestyle{fancy}
\fancyhf{}
\renewcommand{\headrulewidth}{0.4pt}
\setlength{\headheight}{15pt}

% \entete{Matière}{Titre}{Type de document}
\newcommand{\entete}[3]{%
  \fancyhead[L]{\small\textsc{La Revue CAPES} -- #1}%
  \fancyhead[R]{\small #2}%
  \fancyfoot[C]{\thepage}%
  \thispagestyle{plain}%
  \begin{center}
    {\color{revue}\rule{\textwidth}{1.2pt}}\\[4pt]
    {\small\textsc{Concours directs CAP-CEG / CAPES -- Burkina Faso}}\\[6pt]
    {\LARGE\bfseries\color{revue} #2}\\[6pt]
    {\large\itshape #1 -- #3}\\[2pt]
    {\color{revue}\rule{\textwidth}{1.2pt}}
  \end{center}
  \vspace{0.5cm}%
}

% Encadré pour les remarques ajoutées dans les corrigés
\newcommand{\remarque}[1]{\par\medskip\noindent\fcolorbox{revue}{revue!6}{\parbox{\dimexpr\linewidth-2\fboxsep-2\fboxrule}{\textbf{\color{revue}Remarque.} #1}}\par\medskip}


\begin{document}
\entete{Culture générale}{Corrigé -- Session 2021}{Correction}
\begin{spacing}{1.5}

\textbf{Introduction}

La jeunesse est souvent qualifiée de $\ll$force  vive$\gg$ d'une société, représentant son potentiel de transformation et d'innovation. Cependant, en tant que l'héritière du monde actuel, elle porte une responsabilité immense, car ses actions ou son inaction auront des répercussions profondes sur l'avenir. Face à cette réalité, l'éducation joue un rôle fondamental pour guider la jeunesse dans la prise de décisions éclairées. Nous examinerons d'abord pourquoi la jeunesse est responsable de ses choix, avant de montrer comment l'éducation peut l'aider à assumer cette responsabilité.


\textbf{Développement}

Première partie : La jeunesse est responsable de son action comme de son inaction

1. L'action de la jeunesse façonne l'avenir d'une société
  
Les jeunes, par leurs idées et leurs initiatives, influencent directement l'évolution des communautés. Par exemple, Malala Yousafzai, militante pour l'éducation des filles, a démontré qu'une jeune personne peut transformer les mentalités et influencer des politiques mondiales. Une jeunesse engagée dans des actions positives contribue à la construction d'un monde meilleur.

2. L'inaction a des conséquences graves  

À l'inverse, l'inaction des jeunes face à des enjeux majeurs, comme le changement climatique, peut aggraver les crises existantes. Par exemple, si les jeunes restent passifs face à la déforestation en Afrique, les écosystèmes se dégradent, ce qui compromet leur propre avenir. Leur silence ou leur indifférence est une forme de responsabilité qui les rattrapera.

Deuxième partie : Le rôle de l'éducation pour orienter la jeunesse vers le bon choix

1. L'éducation donne les outils pour comprendre les enjeux.
  
Une éducation de qualité permet aux jeunes d'acquérir les compétences nécessaires pour analyser les défis contemporains. Par exemple, dans les écoles qui intègrent l'éducation environnementale, les élèves deviennent souvent des militants actifs pour des causes écologiques. L'éducation leur ouvre ainsi les yeux sur les problématiques complexes de leur époque.

2. L'éducation développe un sens des responsabilités et des valeurs.
  
Au-delà des connaissances, l'éducation forme les individus à des valeurs comme l'éthique, la solidarité et la citoyenneté. Par exemple, les programmes qui mettent l'accent sur le volontariat, comme les "services civiques", incitent les jeunes à s'impliquer dans leur communauté. Une jeunesse éduquée est donc mieux préparée à faire des choix responsables et durables.

\textbf{Conclusion}

En somme, la jeunesse porte la responsabilité de l'avenir à travers ses actions et ses choix. Si elle s'implique activement, elle peut bâtir un monde meilleur, mais son inaction pourrait avoir des conséquences désastreuses. L'éducation apparaît ainsi comme un levier indispensable pour éclairer cette jeunesse et la préparer à prendre des décisions éclairées, guidées par des valeurs et une vision claire des enjeux mondiaux. Dès lors, investir dans l'éducation, c'est investir dans un avenir prometteur.


\quad
\quad

\end{spacing}
\end{document}
